\honest{Planning Fallacy}{Eliezer Yudkowsky}{2007}

Denver International Airport opened 16 months late and \$2 billion over budget. The Eurofighter Typhoon was 54 months late and cost \$19 billion instead of \$7 billion. The Sydney Opera House was supposed to be finished in 1963 for \$7 million, and was finished in 1973 for \$102 million.

Are these just the disasters we happen to hear about, or the work of bureaucracy and bad incentives? ``Yes, very probably,'' I answer. So, by my own account, the three cases that open my post very probably have other causes too. They cannot on their own show the bias I am about to describe. \nb{For Denver, the US General Accounting Office reported in 1995 that most of the cost growth came from changes in the airport's scope.}

The bias shows up in experiments. Students asked for dates by which they were 50\%, 75\% and 99\% sure to finish a project met them only 13\%, 19\% and 45\% of the time. Another study found that people's ``realistic'' scenarios were the same as their best-case ones. Asked for a realistic plan, people picture everything going as planned, with no unexpected delays. Reality ``usually delivers results somewhat worse'' than the worst case. I give no source for that, yet a few paragraphs later I write that ``we saw'' that picturing the worst case is not enough. \nb{The 99\% forecasts come close to showing this. Direct evidence was in the papers the post cites. Newby-Clark and colleagues (2000) found that pessimistic scenarios did not change people's predictions. Buehler, Griffin and Ross (1994) found that about half of one class of students finished after the date they gave for the worst case. The post reports neither finding.}

Planning as we usually do it is the ``inside view'': work out where, when and how, estimate time and resources, and picture the steps from start to finish. It leaves out the delays no one foresees. The cure is the ``outside view.'' Ignore what is special about your project and ask how long broadly similar projects took. This feels wrong, since the inside view has more detail. But in experiments, more detailed plans were more optimistic and less accurate. It will not work for projects like Denver's, but it works for personal planning. Shoppers asked to plan their Christmas shopping in detail expected to finish more than a week before Christmas; shoppers simply asked when they would finish said four days before. Both groups finished three days before. Japanese students expected to finish their essays ten days early, finished one day early, and said that one day early was what they usually did. Experienced outsiders are ``often much less optimistic and much more accurate.'' I give no source. \nb{The one test of outside forecasters in the 1994 paper used fellow students. They were more pessimistic than the students doing the work, and ``no more accurate.''}

My conclusion: ask how long similar projects have taken, or ask an outsider. The answer will sound ``hideously long.'' ``This answer is true. Deal with it.'' \nb{In the same paper, students asked to relate their past to the task at hand lost their optimistic lean, but their forecasts were ``no more accurate'' than the others'. The outside view corrects the average error. It does not give the true answer for a given project.}
